\section*{Abstract}
In this thesis, a framework for robot simulation is developed, which is used for teaching at the University of Applied Sciences Dortmund.  
The framework enables the teach-in of joint configurations as well as the animation and control of industrial robots with serial kinematics,  
which are calculated using forward and inverse transformations.  
Arbitrary robot models can be integrated into the framework via Unified Robot Description Format (URDF) descriptions, including robots with seven or more rotary joints,  
as the Newton-Raphson method is employed for inverse kinematics.  
The Denavit-Hartenberg convention serves as the mathematical model, and its coordinate systems are visually represented.  
Furthermore, the framework supports the simulation and animation of mobile robots.  
Interactive controls, such as sliders for rotating cubes, illustrate kinematic movements directly in the frontend.  
The goal of the framework is to provide a clear and practical understanding of kinematic concepts in robotics.  
A conducted stress test confirms that the framework runs smoothly and stably under typical system conditions.
